% Independently redrawn replacements for the 26 scanned source diagrams in Section 1.3.
% Each registration returns only a tikzpicture; the original requested width is retained.
\newcommand{\registerfaithful}[2]{\expandafter\def\csname faithful@#1\endcsname{#2}}
\newcommand{\sourcefigure}[2][0.52\linewidth]{%
  \begin{center}\resizebox{#1}{!}{\csname faithful@#2\endcsname}\end{center}}
\newcommand{\sourcefiguretrim}[3][0.52\linewidth]{\sourcefigure[#1]{#2}}
\newcommand{\sourcefigurefixB}{\sourcefigure[0.42\linewidth]{ex-1-10.png}}

\registerfaithful{fig-1-1.png}{%
\begin{tikzpicture}[x=1cm,y=1cm,>=Latex]
  % B is the orthogonal projection of O onto the force line through A.
  \coordinate (O) at (0,0); \coordinate (A) at (5.2,0); \coordinate (B) at (2.148,2.560);
  \draw[Ink,very thick] (O)--(A) node[midway,below] {$r$};
  \draw[TextGray,dashed] (O)--(B) node[midway,above left] {$r_\perp$};
  \draw[Ink,very thick,->] (A)--(B) node[midway,below left] {$F$};
  \draw[TextGray,dashed] (B)--++(-0.92,0.77) (B)--(5.2,2.560) (A)--(5.2,2.560);
  \draw[Ink,thick,->] (A)--(5.2,2.560) node[right] {$F_\perp$};
  \draw[Ink,thick,->] (5.2,2.560)--(B);
  \draw[Ink,thin] ($(B)+(-40:.25)$)--($(B)+(-40:.25)+(50:.25)$)--($(B)+(50:.25)$);
  \fill[Ink] (O) circle (2.2pt); \node[below left] at (O) {$O$}; \node[below right] at (A) {$A$}; \node[above right=8pt] at (B) {$B$};
  \draw[Ink] (4.25,0) arc[start angle=180,end angle=140,radius=0.95]; \node at (4.00,0.42) {$\theta$};
\end{tikzpicture}}

\registerfaithful{ex-1-10.png}{%
\begin{tikzpicture}[x=1cm,y=1cm,>=Latex]
  \draw[Ink,thick] (-3.1,2.295)--(3.1,2.295); \foreach \x in {-2.9,-2.5,...,2.9}{\draw[Ink] (\x,2.295)--++(0.28,0.30);}
  % AC=CD=DB=l=2.65 and AB=2l.
  \coordinate (A) at (-2.65,2.295); \coordinate (B) at (2.65,2.295); \coordinate (C) at (-1.325,0); \coordinate (D) at (1.325,0);
  \draw[Ink,very thick] (A)--(C)--(D)--(B); \draw[Ink,thick] (C)--++(0,-1.00);
  \draw[fill=SoftGray,draw=Ink,thick] (-1.725,-1.70) rectangle (-0.925,-1.00);
  \foreach \p/\lab/\pos in {A/A/above left,B/B/above right,C/C/below left,D/D/below right}{\fill[Ink] (\p) circle (1.7pt);\node[\pos] at (\p) {$\lab$};}
\end{tikzpicture}}

\registerfaithful{ex-1-11.png}{%
\begin{tikzpicture}[x=1cm,y=1cm,>=Latex]
  % The two boards meet at V; the circle is tangent to the 12 and 72 degree rays.
  \coordinate (V) at (-2.60,-1.20);
  \coordinate (C) at ($(V)+(42:{.85/sin(30)})$);
  \coordinate (E) at ($(V)+(12:5.75)$);
  \draw[Ink,thick] (-3.0,-1.20)--(3.25,-1.20);
  \draw[Ink,very thick] (V)--++(72:4.55);
  \draw[Ink,very thick] (V)--(E);
  \draw[Ink,thick] (E)--(E|-V);
  \draw[Ink,thick] (C) circle (0.85);
  \draw[Ink] ($(V)+(0:0.72)$) arc[start angle=0,end angle=12,radius=0.72];
  \draw[TextGray,thin] ($(V)+(6:.74)$)--(-1.30,-1.50);
  \node[below,inner sep=1pt] at (-1.30,-1.50) {$\alpha$};
  % Put the included-angle arc next to the vertex, clear of the ball.
  \draw[Ink] ($(V)+(12:0.52)$) arc[start angle=12,end angle=72,radius=0.52];
  \node[font=\small,inner sep=0pt] at ($(V)+(42:0.65)$) {$\theta$};
\end{tikzpicture}}

\registerfaithful{ex-1-12.png}{%
\begin{tikzpicture}[x=1cm,y=1cm,>=Latex]
  \coordinate (V) at (-3.0,-1.2); \coordinate (C) at ($(V)+(32.5:{.92/sin(19.5)})$);
  % The fixed faces are 13 and 52 degrees above horizontal; the ball is tangent to both.
  \draw[Ink,thick] (V)--(3.4,-1.2); \draw[Ink,very thick] (V)--++(13:6.1); \draw[Ink,very thick] (V)--++(52:5.0);
  \draw[Ink,thick] (C) circle (0.92);
  \draw[Ink] ($(V)+(0:.85)$) arc[start angle=0,end angle=13,radius=.85];
  \draw[TextGray,thin] ($(V)+(6.5:.86)$)--(-1.75,-1.55);
  \node[below,font=\small] at (-1.75,-1.55) {$\theta_1$};
  % Both inclinations start from the same horizontal ray; theta2 is not the wedge angle.
  \draw[Ink] ($(V)+(0:1.15)$) arc[start angle=0,end angle=52,radius=1.15];
  \node[font=\small,inner sep=0pt] at ($(V)+(34:1.42)$) {$\theta_2$};
\end{tikzpicture}}

\registerfaithful{ex-1-13.png}{%
\begin{tikzpicture}[x=1cm,y=1cm]
  % R=1.15, r=R/2.  The two large balls and the small ball are drawn
  % exactly tangent.  O is 3R from each large-ball centre, so the visible
  % string segment from O to the surface has length 2R and its extension
  % passes through the corresponding centre.
  \draw[Ink,thick] (-2.4,2.703)--(2.4,2.703); \foreach \x in {-2.2,-1.8,...,2.2}{\draw[Ink] (\x,2.703)--++(0.25,0.28);}
  \coordinate (O) at (0,2.703);
  \coordinate (A) at (-1.15,-0.55); \coordinate (B) at (1.15,-0.55); \coordinate (C) at (0,0.73574);
  \coordinate (As) at ($(A)!0.333333!(O)$); \coordinate (Bs) at ($(B)!0.333333!(O)$);
  \draw[Ink,thick] (O)--(As) (O)--(Bs);
  \draw[Ink,thick] (A) circle (1.15); \draw[Ink,thick] (B) circle (1.15); \draw[Ink,thick] (C) circle (0.575);
  \fill[Ink] (A) circle (1.8pt) node[below left] {$A$}; \fill[Ink] (B) circle (1.8pt) node[below right] {$B$}; \node at (C) {$C$};
\end{tikzpicture}}

\registerfaithful{ex-1-14.png}{%
\begin{tikzpicture}[x=1cm,y=1cm,>=Latex]
  \draw[Ink,thick] (-1.35,3.35)--(1.35,3.35); \foreach \x in {-1.2,-0.8,...,1.2}{\draw[Ink] (\x,3.35)--++(0.25,0.28);}
  \coordinate (O) at (0,3.35); \coordinate (C) at (-0.75,0.10);
  \coordinate (P) at (0.664,0.421); \coordinate (Q) at (0.70,0.10);
  \draw[TextGray,dashed] (O)--(C); \draw[Ink,thick] (O)--(C);
  \draw[Ink,thick] (C) circle (1.45) node[below left=11pt] {$m_1$};
  % OP is tangent; the rope then follows the surface to Q and leaves vertically.
  \draw[Ink,very thick] (O)--(P);
  \draw[Ink,very thick] (P) arc[start angle=12.8,end angle=0,radius=1.45];
  \draw[Ink,very thick] (Q)--(0.70,-1.45);
  \draw[fill=SoftGray,draw=Ink] (0.35,-1.96) rectangle (1.05,-1.45);
  \node[right=4pt] at (1.05,-1.70) {$m_2$};
  \draw[TextGray,dashed] (O)--++(0,-1.15);
  % theta is the angle between OC and the downward vertical.
  \draw[Ink] ($(O)+(-103:0.70)$) arc[start angle=-103,end angle=-90,radius=0.70];
  \draw[TextGray,thin] ($(O)+(-96.5:.70)$)--(-.72,2.68)--(-1.03,2.68);
  \node[left,font=\small] at (-1.03,2.68) {$\theta$};
  \fill[Ink] (O) circle (1.5pt) node[above=11pt] {$O$};
\end{tikzpicture}}

\registerfaithful{ex-1-15.png}{%
\begin{tikzpicture}[x=1cm,y=1cm,>=Latex]
  % Choose a=2 and r=1.1, satisfying a<2r<2a.  Both balls are tangent
  % to their adjacent wall and exactly tangent to each other.
  \draw[Ink,very thick] (-2.00,-1.55)--(-2.00,3.45) (2.00,-1.55)--(2.00,3.45); \draw[Ink,thick] (-2.50,-1.55)--(2.50,-1.55);
  \foreach \x in {-2.4,-2.0,...,2.4}{\draw[Ink] (\x,-1.55)--++(-0.22,-0.24);}
  \coordinate (O) at (-0.90,-0.45); \coordinate (Op) at (0.90,0.815);
  \draw[Ink,thick] (O) circle (1.10); \draw[Ink,thick] (Op) circle (1.10);
  \fill[Ink] (O) circle (1.6pt) node[below] {$O$}; \fill[Ink] (Op) circle (1.6pt) node[above right] {$O'$};
  \draw[<->,Ink] (-2.00,3.02)--node[above] {$2a$}(2.00,3.02); \node[below=9pt] at (-2.00,-1.55) {$A$}; \node[below=9pt] at (2.00,-1.55) {$B$};
  \draw[Ink,->] (Op)--++(130:1.10);
  \node at ($(Op)+(130:.55)+(220:.25)$) {$r$};
\end{tikzpicture}}

\registerfaithful{ex-1-16.png}{%
\begin{tikzpicture}[x=1cm,y=1cm,>=Latex]
  \draw[Ink,thick] (-3.4,-2.1)--(-3.4,2.8); \foreach \y in {-1.9,-1.5,...,2.7}{\draw[Ink] (-3.4,\y)--++(-0.28,0.22);}
  \coordinate (A) at (-3.4,-1.55); \coordinate (E) at ($(A)+(25:6.60)$); \draw[Ink,very thick] (A)--(E);
  \draw[Ink,very thick] (-3.4,{6.60*sin(25)-1.55})--(E); \draw[->,Ink,thick] (E)--++(180:3.1) node[midway,above] {$F$};
  % r=.55; s1=r(1+sin25)/cos25, so cylinder 1 is tangent to both wall and board.
  \pgfmathsetmacro{\sone}{.55*(1+sin(25))/cos(25)}
  \foreach \i in {0,1,2}{\coordinate (C\i) at ($(A)+(25:{\sone+1.10*\i})+(115:.55)$); \draw[Ink,thick] (C\i) circle (.55);}
  \node[below left=7pt] at (A) {$A$};
  % Attach r to an actual radius, so it cannot be read as an F subscript.
  \draw[Ink,->] (C2)--++(70:.55) node[midway,left] {$r$};
  \node[below] at (1.75,-0.02) {$l$};
  \draw[Ink] ($(A)+(90:.34)$) arc[start angle=90,end angle=25,radius=.34];
  \draw[TextGray,thin] ($(A)+(57.5:.34)$)--(-4.00,-1.20);
  \node[left] at (-4.00,-1.20) {$\theta$};
\end{tikzpicture}}

\registerfaithful{ex-1-17.png}{%
\begin{tikzpicture}[x=1cm,y=1cm]
  \draw[Ink,thick] (-4.2,-1.35)--(4.2,-1.35);
  % Each wall is a 60-degree line; the end tangent differs from it by exactly 30 degrees.
  \draw[Ink,very thick] (-2.20,-1.35)--(-3.63,1.12) (2.20,-1.35)--(3.63,1.12);
  \draw[Ink,thick,domain=-3.63:3.63,samples=90] plot (\x,{.07948*\x*\x+.073});
  \draw[Ink,thick,domain=-3.63:3.63,samples=90] plot (\x,{1.65-.0400*\x*\x});
  \foreach \x in {-3.0,-2.0,-1.0,0,1.0,2.0,3.0}{\draw[Ink] (\x,{1.65-.0400*\x*\x})--(\x,{.07948*\x*\x+.073});}
  \coordinate (Ewall) at (3.63,1.12);
  \draw[Ink] ($(Ewall)+(210:.58)$) arc[start angle=210,end angle=240,radius=.58];
  \node at ($(Ewall)+(225:.82)$) {$30^\circ$};
  \node[above] at (0,1.70) {屋面}; \node[below] at (0,0.073) {钢索}; \node[right] at (3.34,0.45) {斜壁};
\end{tikzpicture}}

\registerfaithful{fig-1-2.png}{%
\begin{tikzpicture}[x=1cm,y=1cm,>=Latex]
  \draw[Ink,thick] (-2.8,-1.15)--(2.8,-1.15); \foreach \x in {-2.6,-2.2,...,2.6}{\draw[Ink] (\x,-1.15)--++(-0.22,-0.22);}
  \draw[Ink,thick] (-0.50,-1.15) rectangle (0.72,-0.05);
  \coordinate (P) at (0.10,-0.05); \draw[->,Ink,thick] (P)--(0.10,1.72) node[right] {$N$}; \draw[->,Ink,thick] (P)--(-1.38,-0.05) node[below] {$f$};
  \draw[->,DeepBlue,very thick] (P)--(-1.38,1.72) node[above left] {$F$}; \draw[TextGray,dashed] (-1.38,1.72)--(0.10,1.72) (P)--(-1.38,-0.05);
  \pgfmathsetmacro{\forceangle}{atan2(1.77,-1.48)}
  \draw[Ink] ($(P)+(90:.92)$) arc[start angle=90,end angle=\forceangle,radius=.92];
  \node at ($(P)+({(90+\forceangle)/2}:1.20)$) {$\alpha$};
\end{tikzpicture}}

\registerfaithful{ex-1-18.png}{%
\begin{tikzpicture}[x=1cm,y=1cm,>=Latex]
  \draw[Ink,thick] (-3,-1.0)--(3,-1.0); \foreach \x in {-2.8,-2.4,...,2.8}{\draw[Ink] (\x,-1)--++(-0.22,-0.22);}
  \draw[Ink,thick,fill=SoftGray] (-1.30,-1) rectangle (0.25,0.25); \node at (-0.55,-0.35) {$M$};
  \draw[->,DeepBlue,very thick] (0.25,0.15)--(2.28,1.45) node[right] {$F$}; \draw[TextGray,dashed] (0.25,0.15)--(2.28,0.15);
  \draw[Ink] (1.05,0.15) arc[start angle=0,end angle=33,radius=0.80]; \node at (1.15,0.48) {$\alpha$};
\end{tikzpicture}}

\registerfaithful{ex-1-19.png}{%
\begin{tikzpicture}[x=1cm,y=1cm]
  \coordinate (O) at (-3.4,0); \coordinate (C) at (0.851,0);
  \draw[Ink,very thick] (O)--++(15:7.1) (O)--++(-15:7.1);
  \draw[Ink,thick] (C) circle (1.10); \fill[Ink] (O) circle (1.8pt) node[above left] {$O$};
  \draw[Ink] ($(O)+(-15:1.45)$) arc[start angle=-15,end angle=15,radius=1.45];
  \node at ($(O)+(0:1.78)$) {$\theta$};
  % The radius to the upper contact is normal to the 15-degree face.
  \draw[Ink,->] (C)--++(105:1.10) node[midway,right] {$r$};
\end{tikzpicture}}

\registerfaithful{ex-1-20.png}{%
\begin{tikzpicture}[x=1cm,y=1cm]
  \draw[Ink,thick] (-3.1,-1.55)--(-3.1,2.35); \foreach \y in {-1.4,-1.0,...,2.2}{\draw[Ink] (-3.1,\y)--++(-0.25,0.22);}
  \coordinate (A) at (-3.1,-0.72); \coordinate (B) at (2.65,-0.72); \coordinate (C) at (-3.1,1.80);
  \draw[Ink,very thick] (A)--(B); \draw[Ink,thick] (C)--(B); \fill[Ink] (0.0,-0.72) circle (1.8pt) node[below] {$P$};
  \node[left=9pt] at (A) {$A$}; \node[right] at (B) {$B$}; \node[left=9pt] at (C) {$C$};
  \draw[<->,Ink] (-3.0,-0.35)--node[above] {$x$}(0.0,-0.35); \draw[Ink] (1.72,-0.72) arc[start angle=180,end angle=156,radius=0.93]; \node at (1.55,-0.40) {$\theta$};
\end{tikzpicture}}

\registerfaithful{ex-1-21.png}{%
\begin{tikzpicture}[x=1cm,y=1cm,>=Latex]
  \draw[Ink,thick] (-3.1,-1.55)--(3.1,-1.55); \foreach \x in {-2.9,-2.5,...,2.9}{\draw[Ink] (\x,-1.55)--++(-0.22,-0.22);}
  \draw[Ink,very thick] (-2.45,-1.55)--(2.15,0.70); \draw[->,DeepBlue,very thick] (2.15,0.70)--(1.40,2.23) node[above left] {$F$};
  \draw[Ink] (-1.45,-1.55) arc[start angle=0,end angle=26,radius=1.0]; \node at (-1.16,-1.30) {$\theta$};
\end{tikzpicture}}

\registerfaithful{ex-1-22.png}{%
\begin{tikzpicture}[x=1cm,y=1cm]
  \draw[Ink,thick] (-3.4,1.10)--(3.4,1.10); \fill[Ink] (-2.60,1.10) circle (2.4pt) node[above] {$A$} node[below left] {$m$}; \fill[Ink] (2.60,1.10) circle (2.4pt) node[above] {$B$} node[below right] {$m$};
  \coordinate (C) at (0,-0.55); \draw[Ink,thick] (-2.60,1.10)--(C)--(2.60,1.10); \draw[Ink,thick] (C)--(0,-1.35);
  \draw[fill=SoftGray,draw=Ink] (-0.43,-1.90) rectangle (0.43,-1.35) node[midway] {$M$};
  \draw[<->,Ink] (-2.60,1.52)--node[above] {$x$}(2.60,1.52);
\end{tikzpicture}}

\registerfaithful{fig-1-3.png}{%
\begin{tikzpicture}[x=1cm,y=1cm,>=Latex]
  \begin{scope}[shift={(-4.1,0)}]
    \node[font=\small] at (0,3.00) {(a) 竖直绳};
    \draw[Ink,thick] (-0.55,2.0)--(0.55,2.0); \draw[Ink,very thick] (0,2.0)--(0,-2.0); \node[above] at (0,2.05) {$O$};
    % Separate the position coordinate, total length and rope mass so none
    % of the three labels shares the same line or label box.
    \draw[<->,Ink] (-0.62,2.0)--node[midway,left] {$x$}(-0.62,0.35);
    \draw[TextGray,dashed] (-0.62,0.35)--(0,0.35);
    \fill[Ink] (0,0.35) circle (1.5pt);
    \draw[<->,Ink] (-1.30,2.0)--node[midway,left] {$L$}(-1.30,-2.0);
    \node[right] at (0.12,0) {$m$};
  \end{scope}
  \begin{scope}[shift={(2.2,0)}]
    \node[font=\small] at (0,3.00) {(b) 弯绳微元};
    \draw[TextGray,thick,domain=-2.2:1.65,samples=80] plot (\x,{.18*(\x+2.2)^2-1.35});
    \draw[Ink,very thick,domain=-1.55:.70,samples=60] plot (\x,{.18*(\x+2.2)^2-1.35});
    \coordinate (Aseg) at (-1.55,-1.274); \coordinate (Bseg) at (.70,.164);
    \draw[->,Ink,thick] (Aseg)--++(193.2:1.15) node[below left] {$T_A$};
    \draw[->,Ink,thick] (Bseg)--++(46.2:1.25) node[above right] {$T_B$};
    % Place the differential element's weight on the highlighted rope,
    % rather than on the straight chord joining its endpoints.
    \coordinate (Gseg) at (-.425,-.783);
    \draw[->,DeepBlue,thick] (Gseg)--++(0,-1.15) node[right] {$\Delta mg$};
    \fill[Ink] (Aseg) circle (1.6pt) node[above left] {$A$};
    \fill[Ink] (Bseg) circle (1.6pt) node[below right] {$B$};
  \end{scope}
\end{tikzpicture}}

\registerfaithful{ex-1-23.png}{%
\begin{tikzpicture}[x=1cm,y=1cm]
  \draw[TextGray,dashed] (0,3.2)--(0,-2.35); \draw[Ink,thick] (0,3.2)--(-1.75,-1.70) (0,3.2)--(1.75,-1.70);
  \draw[TextGray,dashed] (1.75,-1.70) arc[start angle=0,end angle=180,x radius=1.75,y radius=0.42];
  \draw[Ink,thick] (-1.75,-1.70) arc[start angle=180,end angle=360,x radius=1.75,y radius=0.42];
  \draw[TextGray,dashed,thick] (1.054,0.25) arc[start angle=0,end angle=180,x radius=1.054,y radius=0.27];
  \draw[DeepBlue,very thick] (-1.054,0.25) arc[start angle=180,end angle=360,x radius=1.054,y radius=0.27];
  \pgfmathsetmacro{\coneangle}{atan(1.75/4.90)}
  \draw[Ink] (0,2.15) arc[start angle=-90,end angle={-90+\coneangle},radius=1.05];
  \draw[TextGray,thin] ({1.05*sin(\coneangle/2)},{3.2-1.05*cos(\coneangle/2)})--(1.40,2.40)--(1.65,2.40);
  \node[right,font=\small] at (1.65,2.40) {$\delta=\alpha/2$};
\end{tikzpicture}}

\registerfaithful{ex-1-24.png}{%
\begin{tikzpicture}[x=1cm,y=1cm,>=Latex]
  \draw[Ink,thick] (-3.4,-1.3)--(3.4,-1.3); \foreach \x in {-3.2,-2.8,...,3.2}{\draw[Ink] (\x,-1.3)--++(-0.22,-0.22);}
  \draw[Ink,thick] (-2.25,-1.3) arc[start angle=180,end angle=0,radius=2.25];
  \draw[DeepBlue,line width=1.5pt] (0,0.95) arc[start angle=90,end angle=0,radius=2.25];
  \draw[->,DeepBlue,very thick] (0,0.95)--(-1.85,0.95) node[left] {$F$};
  \coordinate (O) at (0,-1.3); \coordinate (Q) at ($(O)+(135:2.25)$);
  \draw[Ink,->] (O)--(Q);
  \node at ($(O)!.5!(Q)+(225:.30)$) {$R$};
  \fill[Ink] (O) circle (1.8pt) node[below=9pt] {$O$};
\end{tikzpicture}}

\registerfaithful{ex-1-25.png}{%
\begin{tikzpicture}[x=1cm,y=1cm,>=Latex]
  \draw[Ink,thick] (0,0.75) circle (1.65); \fill[Ink] (0,0.75) circle (2pt);
  \draw[Ink,very thick] (-1.65,0.75)--(-1.65,-2.55) (1.65,0.75)--(1.65,-2.55);
  \coordinate (O) at (0,.75); \coordinate (Q) at ($(O)+(135:1.65)$);
  \draw[Ink,->] (O)--(Q); \node at ($(O)!.5!(Q)+(225:.26)$) {$R$};
\end{tikzpicture}}

\registerfaithful{ex-1-26.png}{%
\begin{tikzpicture}[x=1cm,y=1cm]
  \draw[Ink,thick] (-3.0,-2.40)--(3.0,-2.40);
  \foreach \x in {-2.8,-2.4,...,2.8}{\draw[Ink] (\x,-2.40)--++(-.22,-.22);}
  \draw[Ink,thick] (0,0) circle (2.40);
  % The centre line is a construction axis, not a physical rod.
  \draw[TextGray,dashed] (0,2.65)--(0,-2.40);
  % Back and front halves of the horizontal ring must both be visible.
  \draw[TextGray,dashed,thick] (1.697,1.697) arc[start angle=0,end angle=180,x radius=1.697,y radius=0.38];
  \draw[DeepBlue,very thick] (-1.697,1.697) arc[start angle=180,end angle=360,x radius=1.697,y radius=0.38];
  \coordinate (O) at (0,0); \coordinate (Q) at (155:2.40);
  \draw[Ink,->] (O)--(Q); \node at ($(O)!.5!(Q)+(245:.28)$) {$R$};
  \fill[Ink] (O) circle (2pt) node[below right] {$O$};
  \draw[<->,Ink] (0,1.697)--node[above] {$b$}(1.697,1.697);
\end{tikzpicture}}

\registerfaithful{ex-1-27.png}{%
\begin{tikzpicture}[x=1cm,y=1cm]
  \draw[Ink,thick] (-3.4,1.25)--(3.4,1.25); \node[above] at (-3.0,1.25) {$A$}; \node[above] at (3.0,1.25) {$B$};
  \draw[Ink,very thick,domain=-3:3,samples=101] plot (\x,{1.25+3.4*(cosh(\x/3.4)-cosh(3/3.4))});
  \node[above] at (0,1.25) {$d$};
\end{tikzpicture}}

\registerfaithful{fig-1-4.png}{%
\begin{tikzpicture}[x=1cm,y=1cm]
  \begin{scope}[shift={(-4.5,0)}] \draw[Ink,thick] (-0.7,2.2)--(0.7,2.2); \foreach \x in {-0.6,-0.2,0.2,0.6}{\draw[Ink] (\x,2.2)--++(0.22,0.25);} \fill[SoftGray,draw=Ink] (0,1.95) circle (0.15); \draw[Ink,very thick] (0,1.80)--(0,-1.75); \end{scope}
  \begin{scope}[shift={(-1.5,0)}] \draw[Ink,thick] (-0.7,-1.75)--(0.7,-1.75); \foreach \x in {-0.6,-0.2,0.2,0.6}{\draw[Ink] (\x,-1.75)--++(-0.22,-0.25);} \fill[SoftGray,draw=Ink] (0,-1.55) circle (0.15); \draw[Ink,very thick] (0,-1.40)--(0,2.15); \end{scope}
  \begin{scope}[shift={(1.5,0)}] \draw[Ink,thick] (-0.7,2.2)--(0.7,2.2); \foreach \x in {-0.6,-0.2,0.2,0.6}{\draw[Ink] (\x,2.2)--++(0.22,0.25);} \fill[SoftGray,draw=Ink] (0,1.95) circle (0.15); \draw[TextGray,dashed] (0,1.95)--(0,-0.7); \draw[Ink,very thick,rotate around={16:(0,1.95)}] (0,1.80)--(0,-1.75); \end{scope}
  \begin{scope}[shift={(4.5,0)}] \draw[Ink,thick] (-0.7,-1.75)--(0.7,-1.75); \foreach \x in {-0.6,-0.2,0.2,0.6}{\draw[Ink] (\x,-1.75)--++(-0.22,-0.25);} \fill[SoftGray,draw=Ink] (0,-1.55) circle (0.15); \draw[TextGray,dashed] (0,-1.55)--(0,0.85); \draw[Ink,very thick,rotate around={-16:(0,-1.55)}] (0,-1.40)--(0,2.15); \end{scope}
\end{tikzpicture}}

\registerfaithful{ex-1-28.png}{%
\begin{tikzpicture}[x=1cm,y=1cm,decoration={coil,aspect=0.45,segment length=4pt,amplitude=3pt}]
  \draw[Ink,very thick] (0,-1.75)--(0,1.50); \fill[Ink] (0,-1.75) circle (2.2pt); \node[right] at (0,0) {$m,L$};
  \draw[Ink,thick] (-3.9,2.10)--(-3.9,0.95) (3.9,2.10)--(3.9,0.95); \foreach \y in {1.05,1.35,...,2.05}{\draw[Ink] (-3.9,\y)--++(-0.25,0.20);\draw[Ink] (3.9,\y)--++(0.25,0.20);}
  \draw[Ink,thick,decorate] (-3.9,1.50)--(0,1.50); \draw[Ink,thick,decorate] (0,1.50)--(3.9,1.50); \node[above] at (-2.0,1.72) {$k$}; \node[above] at (2.0,1.72) {$k$};
  \draw[Ink,thick] (-0.85,-1.75)--(0.85,-1.75); \foreach \x in {-0.75,-0.35,0.05,0.45}{\draw[Ink] (\x,-1.75)--++(-0.22,-0.22);}
\end{tikzpicture}}

\registerfaithful{ex-1-29.png}{%
\begin{tikzpicture}[x=1cm,y=1cm,decoration={coil,aspect=0.45,segment length=4pt,amplitude=3pt}]
  \coordinate (T) at (0,2.55); \coordinate (A) at (-2.25,-1.45); \coordinate (B) at (2.25,-1.45);
  \draw[Ink,very thick] (A)--(T)--(B);
  \filldraw[fill=SoftGray,draw=Ink] (T) circle (0.18);
  \node[above=4pt,text=Ink,inner sep=0pt] at (T) {$M$};
  \draw[Ink,thick,decorate] (A)--(B); \node[above=5pt] at (0,-1.45) {$k$}; \node[left=5pt] at (-1.15,0.45) {$L$}; \node[right=5pt] at (1.15,0.45) {$L$};
  % The lower ends are sliders on a smooth horizontal support.  Drawing
  % the support here keeps the problem figure and the solution model equal.
  \foreach \X in {-2.37,-2.13,2.13,2.37}{\draw[Ink] (\X,-1.64) circle (0.08);}
  \draw[Ink,thick] (-2.70,-1.76)--(2.70,-1.76);
  \foreach \x in {-2.55,-2.15,...,2.55}{\draw[Ink] (\x,-1.76)--++(-0.18,-0.18);}
  \draw[TextGray,dashed] (T)--(0,-0.10); \draw[Ink] (0,1.55) arc[start angle=-90,end angle=-60,radius=1.0]; \node at (0.36,1.38) {$\theta$}; \node[left] at (A) {$A$}; \node[right] at (B) {$B$};
\end{tikzpicture}}

\registerfaithful{ex-1-30.png}{%
\begin{tikzpicture}[x=1cm,y=1cm,>=Latex]
  \draw[axis] (-0.25,0)--(4.65,0) node[right] {$x$}; \draw[axis] (0,-0.25)--(0,4.60) node[above] {$y$};
  \def\rodlength{4.4}\def\xp{3.20}
  \pgfmathsetmacro{\yp}{(\rodlength-sqrt(\rodlength^2-\xp^2))/2}
  \draw[Ink,very thick,domain=0:4.38,samples=101] plot (\x,{(\rodlength-sqrt(\rodlength^2-\x^2))/2}) node[right] {$y(x)$};
  \coordinate (P) at (\xp,\yp); \coordinate (A) at (0,{\rodlength-\yp});
  \draw[Ink,very thick] (A)--(P) node[midway,above right=4pt] {$L$};
\end{tikzpicture}}

\newcommand{\faithfulExOneThirtyOne}{%
\begin{tikzpicture}[x=0.72cm,y=0.72cm,>=Latex]
  % The problem figure shows the reference equilibrium position.  The
  % displaced rolling state is reserved for the solution diagram.
  \def\Rbow{3.05}\def\rball{0.86}
  \pgfmathsetmacro{\aradius}{\Rbow-\rball}
  \coordinate (C) at (0,1.15);
  \coordinate (O) at ($(C)+(-90:\aradius)$);
  \coordinate (K) at ($(C)+(-90:\Rbow)$);
  % At the reference equilibrium position G lies vertically above O.
  \coordinate (G) at ($(O)+(90:0.54)$);
  \draw[Ink,very thick] ($(C)+(180:\Rbow)$) arc[start angle=180,end angle=360,radius=\Rbow];
  \draw[Ink,thick] (O) circle (\rball);
  \fill[Ink] (O) circle (1.8pt) node[below left=4pt] {$O$};
  \fill[Ink] (G) circle (1.8pt);
  \draw[TextGray,thin] (G)--($(O)+(-1.25,1.10)$);
  \node[left,inner sep=1pt] at ($(O)+(-1.25,1.10)$) {$G$};
  \draw[TextGray,dashed] (O)--($(O)+(1.18,0)$);
  \draw[TextGray,dashed] (G)--($(G)+(1.18,0)$);
  \draw[<->,Ink] ($(O)+(1.18,0)$)--($(G)+(1.18,0)$)
    node[midway,right=2pt,inner sep=0pt] {$h$};
  \coordinate (Qr) at ($(O)+(135:\rball)$);
  \draw[Ink,->] (O)--(Qr);
  \node[font=\small] at ($(O)!.53!(Qr)+(225:.23)$) {$r$};
  \fill[Ink] (K) circle (1.3pt);
  \draw[Ink,->] (C)--++(180:\Rbow) node[midway,above] {$R$};
  \fill[Ink] (C) circle (1.8pt) node[above right=2pt] {$C$};
\end{tikzpicture}}
